\documentclass[addpoints]{exam}
\usepackage{amsmath}
\usepackage{amssymb}
\usepackage{color}
\usepackage{siunitx}

% \printanswers

\newcommand{\event}{Trigonometry \& Geometry}
\newcommand{\tournament}{}
\def\type{0}

\setlength\fillinlinelength{\textwidth}
\CorrectChoiceEmphasis{\color{red}\bfseries}
\SolutionEmphasis{\color{red}\bfseries}

\setlength\answerclearance{2pt}
\pagestyle{headandfoot}
\runningheadrule
\runningheader{\event}{\tournament}{\ifnum\type=1 Class Set Test \else Full Name:\hspace{1cm} \fi}
\runningfootrule
\runningfooter{}{Page \thepage\ of \numpages}{}

\begin{document}
\begin{center}
    {\huge Grade 12 Foundations for College Math \\ \event
    \ifprintanswers \space Key
    \else \space \ifnum\type<2 Test \fi \ifnum\type=1 \textbf{(CLASS SET)} \fi
          \ifnum\type=2 Answer Sheet \fi
    \fi \\ \vspace{3mm}\tournament\vspace{1mm}}
\end{center}

\vspace{.25\textheight}

\begin{center}
    First Name: \ifprintanswers
    \underline{\hspace{3.5cm}}\textbf{KEY}\underline{\hspace{5.5cm}}\\\ \vspace{5mm}
    \else \underline{\hspace{10cm}}\\ \vspace{5mm} \fi
    Last Name: \underline{\hspace{10cm}}\\ \vspace{5mm}
\end{center}

\noindent\textbf{\underline{Directions:}}
\begin{itemize}
    \ifnum\type=1 \item \textbf{This test is a class set.} Do not write on this paper. \fi
    \item Show all work for full marks. Round dollar amounts to the nearest cent.
    \item The test is designed to be completed in 75 minutes.
\end{itemize}
\begin{center}
\textit{For grading use only}\\
\multirowgradetable{1}[pages]
\end{center}

\newpage
\begin{questions}

\section*{Part A --- Multiple Choice (15 marks)}
\vspace{5mm}

% Q1
\question[1]
The Sine Law states that in any triangle $ABC$:
\begin{choices}
    \correctchoice $\dfrac{a}{\sin A} = \dfrac{b}{\sin B} = \dfrac{c}{\sin C}$
    \choice $a^2 = b^2 + c^2 - 2bc\cos A$
    \choice $\dfrac{\sin A}{a} = \dfrac{\cos B}{b}$
    \choice $a + b + c = 180°$
\end{choices}
\vspace{3mm}

% Q2
\question[1]
The Cosine Law for side $c$ in triangle $ABC$ is:
\begin{choices}
    \choice $c = a\cos B + b\cos A$
    \choice $c^2 = a^2 + b^2$
    \correctchoice $c^2 = a^2 + b^2 - 2ab\cos C$
    \choice $c^2 = a^2 - b^2 + 2ab\cos C$
\end{choices}
\vspace{3mm}

% Q3
\question[1]
An angle of elevation is measured:
\begin{choices}
    \correctchoice Upward from the horizontal to the line of sight
    \choice Downward from the horizontal to the line of sight
    \choice From the vertical to the line of sight
    \choice From true north, clockwise
\end{choices}
\vspace{3mm}

% Q4
\question[1]
A surveyor at point $A$ observes the top of a tower at an angle of elevation of \ang{30}. The tower is \SI{50}{m} tall. The horizontal distance from $A$ to the base of the tower is:
\begin{choices}
    \choice \SI{25}{m}
    \choice \SI{43.3}{m}
    \correctchoice \SI{86.6}{m}
    \choice \SI{100}{m}
\end{choices}
\vspace{3mm}

% Q5
\question[1]
A ship travels on a bearing of N\ang{40}E. This means the ship heads:
\begin{choices}
    \choice Due East, then turns North \ang{40}
    \correctchoice \ang{40} east of due north
    \choice \ang{40} north of due east
    \choice In the direction \ang{40} measured counterclockwise from East
\end{choices}
\vspace{3mm}

% Q6
\question[1]
The volume of a cone with radius $r$ and height $h$ is:
\begin{choices}
    \choice $\pi r^2 h$
    \choice $\dfrac{2}{3}\pi r^2 h$
    \correctchoice $\dfrac{1}{3}\pi r^2 h$
    \choice $\dfrac{4}{3}\pi r^3$
\end{choices}
\vspace{3mm}

% Q7
\question[1]
The total surface area of a closed cylinder with radius $r$ and height $h$ is:
\begin{choices}
    \choice $2\pi r h$
    \correctchoice $2\pi r^2 + 2\pi r h$
    \choice $\pi r^2 + 2\pi r h$
    \choice $2\pi r^2 h$
\end{choices}
\vspace{3mm}

% Q8
\question[1]
Two similar solids have a linear scale factor of $2:5$. The ratio of their surface areas is:
\begin{choices}
    \choice $2:5$
    \choice $8:125$
    \correctchoice $4:25$
    \choice $1:10$
\end{choices}
\vspace{3mm}

% Q9
\question[1]
Two similar solids have a linear scale factor of $2:5$. The ratio of their volumes is:
\begin{choices}
    \choice $2:5$
    \choice $4:25$
    \correctchoice $8:125$
    \choice $1:10$
\end{choices}
\vspace{3mm}

% Q10
\question[1]
Euler's formula for a polyhedron relates vertices $V$, edges $E$, and faces $F$ by:
\begin{choices}
    \choice $V + E + F = 1$
    \choice $V - E + F = 0$
    \correctchoice $V - E + F = 2$
    \choice $V + E - F = 2$
\end{choices}
\vspace{3mm}

% Q11
\question[1]
A scale model of a building uses a scale of 1:200. If the model is \SI{15}{cm} tall, the actual building is:
\begin{choices}
    \choice \SI{15}{m}
    \choice \SI{200}{m}
    \correctchoice \SI{30}{m}
    \choice \SI{3000}{m}
\end{choices}
\vspace{3mm}

% Q12
\question[1]
In a triangle with sides $a = 7$, $b = 9$, $c = 5$, which law should you use first to find angle $A$?
\begin{choices}
    \choice Sine Law, because all sides are given
    \correctchoice Cosine Law, because all three sides are known
    \choice Pythagorean Theorem
    \choice Neither; a triangle with these sides has no solution
\end{choices}
\vspace{3mm}

% Q13
\question[1]
The lateral surface area (curved side only) of a cone with base radius $r$ and slant height $\ell$ is:
\begin{choices}
    \choice $\pi r^2$
    \correctchoice $\pi r \ell$
    \choice $2\pi r \ell$
    \choice $\dfrac{1}{3}\pi r^2 \ell$
\end{choices}
\vspace{3mm}

% Q14
\question[1]
A cube has 8 vertices, 12 edges, and 6 faces. Euler's formula gives $V - E + F = $:
\begin{choices}
    \choice 0
    \choice 1
    \correctchoice 2
    \choice 4
\end{choices}
\vspace{3mm}

% Q15
\question[1]
To minimize the surface area of a cylinder of fixed volume, the optimal relationship between radius $r$ and height $h$ is:
\begin{choices}
    \choice $h = r$
    \choice $h = 2r$
    \correctchoice $h = 2r$
    \choice $h = 4r$
\end{choices}
\vspace{3mm}

\section*{Part B --- Short Answer (33 marks)}

\vspace{5mm}

%% Q16 --- 3D trigonometry / mountain height
\question[6]
\textbf{3D Trigonometry --- Mountain Height.}
From point $A$, the angle of elevation to the top of a mountain is \ang{28}. From point $B$, which is \SI{500}{m} closer to the mountain (along the same horizontal line), the angle of elevation to the top is \ang{40}. Points $A$, $B$, and the base of the mountain are all on level ground.

\begin{parts}
\part[3] Use the two angles and the distance $AB = \SI{500}{m}$ to set up a system of equations involving the horizontal distance $x$ from $B$ to the base of the mountain and the mountain height $h$.
\part[3] Solve for $h$, the height of the mountain. Round to the nearest metre.
\end{parts}

\begin{solution}
\textbf{(a)} Let $x$ = horizontal distance from $B$ to the base of the mountain, and $h$ = height of the mountain.

From point $B$: $\tan(\ang{40}) = \dfrac{h}{x}$, so $h = x\tan(\ang{40})$.

From point $A$ (which is $500$ m farther from the mountain): $\tan(\ang{28}) = \dfrac{h}{x + 500}$, so $h = (x+500)\tan(\ang{28})$.

\textbf{(b)} Set the two expressions equal:
\[
x\tan(\ang{40}) = (x+500)\tan(\ang{28})
\]
\[
x(0.8391) = (x+500)(0.5317)
\]
\[
0.8391x = 0.5317x + 265.85
\]
\[
0.3074x = 265.85 \implies x = \frac{265.85}{0.3074} \approx 864.9\ \text{m}
\]
\[
h = 864.9 \times \tan(\ang{40}) \approx 864.9 \times 0.8391 \approx \SI{725.6}{m}
\]
The mountain is approximately \textbf{\SI{726}{m}} tall.
\end{solution}

\vspace{5mm}

%% Q17 --- Navigation / bearings
\question[6]
\textbf{Navigation.}
A ship leaves port at $O$. It sails on a bearing of N\ang{35}E for \SI{50}{km} to point $A$. It then turns and sails on a bearing of S\ang{60}E for \SI{80}{km} to point $B$.

\begin{parts}
\part[3] Find the straight-line distance $OB$ from port to the ship's final position.
\part[3] Find the bearing from $B$ back to $O$ (i.e., the direction the ship must sail to return to port).
\end{parts}

\begin{solution}
\textbf{(a)} Set up a coordinate system with $O$ at the origin, east as positive $x$, and north as positive $y$.

Leg $OA$: bearing N\ang{35}E means \ang{35} east of north.
\[
A_x = 50\sin(\ang{35}) \approx 50(0.5736) = 28.68\ \text{km}
\]
\[
A_y = 50\cos(\ang{35}) \approx 50(0.8192) = 40.96\ \text{km}
\]

Leg $AB$: bearing S\ang{60}E means \ang{60} east of south.
\[
B_x = A_x + 80\sin(\ang{60}) \approx 28.68 + 80(0.8660) = 28.68 + 69.28 = 97.96\ \text{km}
\]
\[
B_y = A_y - 80\cos(\ang{60}) \approx 40.96 - 80(0.5000) = 40.96 - 40.00 = 0.96\ \text{km}
\]

Distance $OB$:
\[
OB = \sqrt{(97.96)^2 + (0.96)^2} \approx \sqrt{9596.2 + 0.9} \approx \sqrt{9597.1} \approx \SI{97.97}{km}
\]

\textbf{(b)} The vector from $B$ to $O$ is $(-97.96,\ -0.96)$ (pointing west and slightly south).

The angle south of west:
\[
\theta = \arctan\!\left(\frac{0.96}{97.96}\right) \approx \arctan(0.0098) \approx \ang{0.56}
\]
The bearing from $B$ to $O$ is approximately \textbf{S\ang{89.4}W} (essentially due west, very slightly south of west).
\end{solution}

\vspace{5mm}

%% Q18 --- Composite solid (cone + cylinder)
\question[5]
\textbf{Composite Solid.}
A storage container consists of a cylinder topped with a cone. The cylinder has radius \SI{3}{m} and height \SI{8}{m}. The cone has the same base radius and a height of \SI{4}{m}.

\begin{parts}
\part[1] Find the slant height of the cone.
\part[2] Find the total surface area of the composite solid (bottom of cylinder + lateral cylinder + lateral cone; the shared base between cone and cylinder is internal and not counted).
\part[2] Find the total volume of the composite solid.
\end{parts}

\begin{solution}
\textbf{(a)} Slant height of cone:
\[
\ell = \sqrt{r^2 + h_{\text{cone}}^2} = \sqrt{3^2 + 4^2} = \sqrt{9+16} = \sqrt{25} = \SI{5}{m}
\]

\textbf{(b)} Total surface area = (circular bottom of cylinder) + (lateral surface of cylinder) + (lateral surface of cone):
\[
SA = \pi r^2 + 2\pi r h_{\text{cyl}} + \pi r \ell
   = \pi(3)^2 + 2\pi(3)(8) + \pi(3)(5)
   = 9\pi + 48\pi + 15\pi
   = 72\pi \approx \SI{226.2}{m^2}
\]

\textbf{(c)} Total volume:
\[
V = \pi r^2 h_{\text{cyl}} + \tfrac{1}{3}\pi r^2 h_{\text{cone}}
  = \pi(9)(8) + \tfrac{1}{3}\pi(9)(4)
  = 72\pi + 12\pi
  = 84\pi \approx \SI{263.9}{m^3}
\]
\end{solution}

\vspace{5mm}

%% Q19 --- Optimization (cylindrical tank)
\question[8]
\textbf{Volume Optimization.}
A cylindrical storage tank (closed top and bottom) must hold exactly \SI{10}{m^3}. The goal is to find the dimensions (radius $r$ and height $h$) that minimize the total surface area.

\begin{parts}
\part[2] Write $h$ as a function of $r$ using the volume constraint $V = \pi r^2 h = 10$.
\part[2] Write the total surface area $SA$ as a function of $r$ alone (substitute the expression for $h$).
\part[4] Complete the table of values below and identify which radius gives the minimum surface area. Round to two decimal places.

\begin{center}
\begin{tabular}{@{}ccc@{}}
\toprule
$r$ (m) & $h$ (m) & $SA$ (\si{m^2}) \\
\midrule
0.5 & & \\
1.0 & & \\
1.2 & & \\
1.5 & & \\
2.0 & & \\
2.5 & & \\
\bottomrule
\end{tabular}
\end{center}
\end{parts}

\begin{solution}
\textbf{(a)} From $\pi r^2 h = 10$:
\[
h = \frac{10}{\pi r^2}
\]

\textbf{(b)}
\[
SA = 2\pi r^2 + 2\pi r h = 2\pi r^2 + 2\pi r \cdot \frac{10}{\pi r^2} = 2\pi r^2 + \frac{20}{r}
\]

\textbf{(c)} Calculated values ($SA = 2\pi r^2 + 20/r$):

\begin{center}
\begin{tabular}{@{}ccc@{}}
\toprule
$r$ (m) & $h$ (m) & $SA$ (\si{m^2}) \\
\midrule
0.5 & 12.73 & 41.57 \\
1.0 & 3.18  & 26.28 \\
1.2 & 2.21  & 25.62 \\
1.5 & 1.41  & 27.51 \\
2.0 & 0.80  & 35.13 \\
2.5 & 0.51  & 47.31 \\
\bottomrule
\end{tabular}
\end{center}

The minimum occurs near $r \approx \SI{1.17}{m}$ (testing $r = 1.17$: $h = 10/(\pi \times 1.17^2) \approx 2.33\ \text{m}$, $SA \approx 25.60\ \text{m}^2$). From the table, $r = \SI{1.2}{m}$ gives the smallest value at $SA \approx \SI{25.62}{m^2}$ with $h \approx \SI{2.21}{m} \approx 2r$.

\textit{Note:} The theoretical optimum (from calculus) is $h = 2r$. At $r \approx 1.17\ \text{m}$, $h \approx 2.34\ \text{m} = 2r$, confirming the pattern.
\end{solution}

\vspace{5mm}

%% Q20 --- Similar solids
\question[8]
\textbf{Similar Solids.}
Two similar cones, Cone A (small) and Cone B (large), have a linear scale factor of $3:5$.

\begin{parts}
\part[2] State the ratio of their surface areas (Cone A : Cone B) and the ratio of their volumes (Cone A : Cone B).
\part[2] The total surface area of Cone A is $200\pi\ \si{cm^2}$. Find the total surface area of Cone B.
\part[2] The volume of Cone A is $540\pi\ \si{cm^3}$. Find the volume of Cone B.
\part[2] If Cone A has a base radius of \SI{6}{cm}, find the base radius of Cone B. Then find the height of each cone using the surface area formula (note: total surface area of a cone $= \pi r^2 + \pi r \ell$, where slant height $\ell = \sqrt{r^2+h^2}$).
\end{parts}

\begin{solution}
\textbf{(a)} Linear scale factor $k = 3:5$.\\
Ratio of surface areas $= k^2 = 3^2:5^2 = \mathbf{9:25}$.\\
Ratio of volumes $= k^3 = 3^3:5^3 = \mathbf{27:125}$.

\textbf{(b)}
\[
\frac{SA_A}{SA_B} = \frac{9}{25} \implies SA_B = SA_A \times \frac{25}{9} = 200\pi \times \frac{25}{9} = \frac{5000\pi}{9} \approx 1{,}745.3\ \si{cm^2}
\]

\textbf{(c)}
\[
\frac{V_A}{V_B} = \frac{27}{125} \implies V_B = 540\pi \times \frac{125}{27} = \frac{540 \times 125\pi}{27} = 2{,}500\pi \approx 7{,}854\ \si{cm^3}
\]

\textbf{(d)} Scale factor $3:5$, so if $r_A = 6\ \text{cm}$:
\[
\frac{r_A}{r_B} = \frac{3}{5} \implies r_B = 6 \times \frac{5}{3} = 10\ \text{cm}
\]

For Cone A: $SA_A = \pi r_A^2 + \pi r_A \ell_A = 200\pi$.
\[
\pi(36) + \pi(6)\ell_A = 200\pi \implies 36 + 6\ell_A = 200 \implies \ell_A = \frac{164}{6} \approx 27.33\ \text{cm}
\]
\[
h_A = \sqrt{\ell_A^2 - r_A^2} = \sqrt{27.33^2 - 36} \approx \sqrt{746.9 - 36} \approx \sqrt{710.9} \approx 26.66\ \text{cm}
\]

For Cone B: by the scale factor, $h_B = h_A \times \dfrac{5}{3} \approx 26.66 \times \dfrac{5}{3} \approx \SI{44.4}{cm}$.
\end{solution}

\end{questions}
\end{document}
